\label{inv:A2.1}\label{wb:prop:log-ks-divisor}\label{wb:prop:period-derivative-determinant}
\StatementRecord{inv:A2.1}{proved}{}{the hypotheses, notation and definitions stated here; the explicit constructions and calculations in the proof}

On the marked upper half-plane,
\begin{equation}\label{inv:eq:A-modular-K}
6\mu_z^2-\tau_z\beta_z
=\tau_z^2\frac{E_4(\tau)^2}{E_6(\tau)}.
\end{equation}
The right side extends holomorphically through the apparent poles.
Its exact orders are two at $z_1$ and zero at $z_2$; it has no other
zero away from their translates.  The logarithmic connecting map
$\kappa_0:f_*\mathcal E\to\mathcal L\otimes G$ has divisor
$\operatorname{divisor}(\det\kappa_0)=p_1+p_2$.
\end{proposition}


\begin{proof}
\begin{enumerate}[label=\arabic*.,leftmargin=*]
\item Consider the holomorphic quadratic differential
$\mathscr Q=6(d\mu)^2-d\tau\,d\beta$.  Direct differentiation of
the two affine transformation laws in AC26 Definition 3.1 gives
\[
g_j^*\mathscr Q=\tau^{-2}\mathscr Q\quad(j=1,2).
\]
In particular $g_2^2$ fixes $\tau$ and sends
\[
\mu\longmapsto1-\tau-\mu,\qquad
\beta\longmapsto\beta+6\tau+12\mu-6,
\]
which leaves $\mathscr Q$ exactly unchanged.  The symbolic residuals
of all three identities are zero in \texttt{a\_calculations.py}.
The kernel of
$\Delta=\Z/3*\Z/4\to\Z/3*\Z/2=\operatorname{PSL}_2(\Z)$
is the normal subgroup generated by $g_2^2$.  The squared multiplier
of $\mathscr Q$ is the ordinary modular multiplier
$(c\tau+d)^{-2}$, so it is one on every kernel element.  Two generic
points of the marked cover with the same $\tau$ have the same $t$
and differ by an element of $\Delta$; its modular image fixes a
generic point, hence is the identity.  It follows that
$k=\mathscr Q/(d\tau)^2$ is a single-valued meromorphic function of
$\tau$, transforming as a modular form of weight two.  This kernel
argument is needed: mere equivariance of the separate period functions
would not prove descent.

\item At $z_1$, write a linearizing parameter $r$ with
$g_1r=\zeta r$, $\zeta=e^{-2\pi i/3}$.  If $\mathscr Q=K_rdr^2$
has first nonzero coefficient of order $n$, its transformation law
at the fixed point reads $\zeta^{n+2}=\rho^{-2}=\zeta$.
Thus $n\equiv2\pmod3$, and $K_r$ vanishes to order at least two.
The map $\tau$ is unramified there by AC26 Theorem 3.4(i).  At
$z_2$, $\tau-i$ has order two and $K_s$ is holomorphic; therefore
$k$ has at most a simple pole at $i$.  Away from these points
$\tau$ is unramified.  At the cusp $\mu$ and $\beta+\tau$ are
holomorphic in $q$, so $k=1+O(q)$.

\item The modular function $k E_6/E_4^2$ is therefore holomorphic at
$\rho$, at $i$, and at the cusp, as well as everywhere else.
The usual zero statements for $E_4,E_6$ (simple at $\rho,i$,
respectively) are the modular facts used in AC26 Section 3.1.
This function descends to the compact modular curve $X(1)=\Pone$,
so it is constant.  Its cusp value is one.  This proves
\eqref{inv:eq:A-modular-K} and the exact two local orders, without the
workbench's degree count for the coarse determinant line.

\item For completeness we translate to the actual coarse lattices.  Choose
cover eigenframes $e_0,e_1$ of $\mathcal E$ with characters
$1,\zeta^a$, and $b_0,b_1$ of $G$ with characters
$1,\zeta^{-a}$.  Their invariant bases are
$e_0,s^{m-a}e_1$ and $b_0,s^ab_1$.  The logarithmic base frame is
$\ell=dt_j/t_j=m\,ds/s$.  Changing determinants from these frames
to the marked ones gives
\[
\det\kappa_0=\epsilon(t_j)m^{-2}s^{m+2-2a}K_s,
\qquad \epsilon(0)\ne0.
\]
This equality follows by taking the two factors $s/m$ from replacing
$ds$ by $\ell$, the factor $s^{m-a}$ in the source, and the factor
$s^{-a}$ in the target.  Thus at $(m,a)=(3,2)$ the total order in
$s$ is $1+2=3$, and at $(4,1)$ it is $4+0=4$.
Both coarse zeros are simple.  At the cusp, the frame change has
no elementary divisor and $dZ/d\tau$ tends to
$\left(\begin{smallmatrix}0&1\\-1&0\end{smallmatrix}\right)$,
with determinant one.  Formula \eqref{inv:eq:A-modular-K} excludes any
further zero, proving the asserted divisor.
\end{enumerate}
\end{proof}